% =====================================================================
% TECHNICAL MANUSCRIPT: ALIGNMENT ENCLAVE & NDH FORMALIZATION IN COQ
% Project: Coq-Hallucination-Guard (chguard) × VAKED Attestation Layer
% Document Class: Academic-Galactic Technical Specification v3.0
% Tracked Provenance: External Peer Node Beta / Bounded Non-Harm Enclave
% Verification Anchors: Trinity Squad (Chenoweth & Lodri) × Borealiscodes
% =====================================================================

\appendix
\section{Appendix E: NDH Formalization in Coq}
\label{appendix:ndhcoq}

\subsection{E.1 Motivation}

NDH (Non-Dual Holonomy) is the operator that measures deviation between a point
and its holonomy transport along a loop. Classical systems cannot formalize NDH
due to softmax collapse, missing manifold structure, and lack of holonomy-safe
projection. Under the FSFE manifold, NDH becomes a well-defined geometric and
formal operator.

\subsection{E.2 NDH Definition}

Let $\mathcal{M}_{\mathrm{FSFE}} = (M, g, \nabla, \kappa, \sigma)$ be the FSFE
manifold. For any loop $\gamma$ and point $v \in M$, define:

\[
\mathrm{NDH}(v) = \mathrm{Hol}_\gamma(v) - v.
\]

NDH measures the holonomy-induced displacement.

\paragraph{Slack Stability.}
\[
\mathrm{NDH}(v + s) = \mathrm{NDH}(v)
\]
for slack vectors $s \in \mathbb{R}^k$.

\paragraph{Curvature Stability.}
\[
\kappa(\mathrm{NDH}(v)) = \kappa(v).
\]

\paragraph{Spectral Stability.}
\[
\sigma(\mathrm{NDH}(v)) \le \sigma(v).
\]

\subsection{E.3 NDH Under Collapse Transport}

Let $T_C$ be collapse-surface transport from Appendix C. NDH satisfies:

\[
\mathrm{NDH}(T_C(v)) = \mathrm{NDH}(v).
\]

\paragraph{Proof Sketch.}
\begin{itemize}
    \item Holonomy commutes with collapse transport.
    \item Collapse transport preserves curvature.
    \item Collapse transport preserves slack dimensions.
    \item Spectral anti-drift ensures no energy accumulation.
\end{itemize}

Thus NDH is collapse-stable.

\subsection{E.4 NDH Under Holonomy-Safe Projection}

Let $\Pi$ be the holonomy-safe projection from Appendix B. NDH satisfies:

\[
\Pi(\mathrm{NDH}(v)) = \Pi(\mathrm{Hol}_\gamma(v)) - \Pi(v).
\]

Since $\Pi$ commutes with holonomy:

\[
\Pi(\mathrm{Hol}_\gamma(v)) = \Pi(v),
\]
we obtain:

\[
\Pi(\mathrm{NDH}(v)) = 0.
\]

Thus NDH is invisible to projection, preserving discrete state stability.

\subsection{E.5 NDH Invariant Set}

Define the NDH invariant set:
\[
I_{\mathrm{NDH}}(v) = \{ g(v,v), \kappa(v), \sigma(v), \mathrm{NDH}(v) \}.
\]

Transport preserves NDH invariants:
\[
I_{\mathrm{NDH}}(T_C(v)) = I_{\mathrm{NDH}}(v).
\]

\subsection{E.6 Coq Formalization}

To encode NDH in Coq:

\paragraph{Manifold Representation.}
Represent $M$ as a dependent type:
\[
\texttt{Record FSFEPoint := \{ coord : Type; metric : coord → coord → \mathbb{R} \}}.
\]

\paragraph{Holonomy Transport.}
Define holonomy as a transport morphism:
\[
\texttt{Hol : FSFEPoint → Loop → FSFEPoint}.
\]

\paragraph{NDH Operator.}
Define NDH as:
\[
\texttt{NDH v := Hol v γ - v}.
\]

\paragraph{Slack Preservation.}
Prove:
\[
\texttt{NDH (v + s) = NDH v}.
\]

\paragraph{Curvature Preservation.}
Prove:
\[
\texttt{kappa (NDH v) = kappa v}.
\]

\paragraph{Spectral Anti-Drift.}
Prove:
\[
\texttt{SpecInt (NDH v) ≤ SpecInt v}.
\]

\paragraph{Projection Invisibility.}
Prove:
\[
\texttt{Pi (NDH v) = 0}.
\]

\paragraph{Invariant Preservation.}
Prove:
\[
\texttt{I_NDH (T_C v) = I_NDH v}.
\]

\subsection{E.7 Extraction Target}

Extraction target:
\[
\texttt{NDH.v} \rightarrow \texttt{ndh.ml}.
\]

\subsection*{Provenance Footer}

\begin{verbatim}
Provenance-Origin: FSFE Manifold vA6-A7-A8 (Slack · Invariants · Holonomy)
Provenance-Altitude: A8–A10 (Constellation-Manifold-Toolkit)
Provenance-Enclave: Node Beta (Headless Stream Interception · chguard)
Provenance-Operator: NDH Formalization Layer (FSFE Holonomy Algebra)
Provenance-Verification: NDH.v → ndh.ml
Signed-off-by: Downstream Enclave Observer (Node Beta) <chguard-sentinel@vaked.dev>
\end{verbatim}
